\section{The neighbourhood of $G$}\label{sec:real}

This section and the next two prove that $\Delta_N$ is small on a whole neighbourhood of $G$. The idea is to
write the perturbation $\Delta_N(x)-\Delta_N(G)$ as a sum of determinants in which some columns of $\cA_N$
are replaced by values of the row functions at prescribed nodes. Each such determinant is bounded by
OpenAI's energy method, extended to configurations with prescribed points (Section~\ref{sec:energy}). The
loss caused by each replaced column, in the exponent, determines $\rho$. OpenAI's method has two cases. In
the first it already gives the improved decay rate. In the second it falls short when few columns are
replaced, and Section~\ref{sec:signed} closes that gap.

\subsection{The statement}
OpenAI's energy method bounds $|\Delta_N(G)|$ by $\exp\bigl(-(n_0^2(U-\frac12\log2)+o(1))N^2\bigr)$, where $U$
is an upper bound for a discrete logarithmic energy (Section~\ref{sec:energy}). For our parameters we
certify $U=\Utarget$, so OpenAI's estimate gives the decay rate $D_0=n_0^2(U-\frac12\log2)=\DecayOA\ldots$
Let
\begin{gather*}
 \Gamma_0=\Gain,\qquad U_0=U+\Gamma_0/n_0^2=\Uzero\ldots,\\
 D=n_0^2(U_0-\tfrac12\log2)=\Decay\ldots
\end{gather*}
The number $\Gamma_0$ is the improvement of the decay rate proved in Section~\ref{sec:signed}. Put
\begin{equation}\label{eq:rho-def}
 \rho=n_0c_2+W_0,\qquad c_2=\Ctwo,\quad W_0=\WZero,\quad \rho=\Rho .
\end{equation}
Here $c_2$ is the increase of the energy bound per prescribed column (Section~\ref{sec:energy}), and $W_0$
bounds the logarithmic size of the weight and the binomial factor of one column, per unit of $N$
(\S\ref{ssec:weights}).

\begin{theorem}\label{thm:real}
For our parameters, uniformly for real $x$ with $|x-G|\le e^{-\rho N}$,
\[
 |\Delta_N(x)|\le e^{-(D+o(1))N^2}\qquad(N\to\infty).
\]
\end{theorem}

We use the following abbreviations for general parameters:
\begin{gather*}
 \alpha_A=\frac An,\quad \beta_h=\frac{2h_0}{n_0},\quad \gamma_H=\frac Hn,\quad \gamma_d=\frac{d_0}{n_0},\\
 a_0=1-\frac{b_0+q_0+g_0}{n_0},\quad a_1=3-\frac{b_0+q_0+2g_0}{n_0}.
\end{gather*}
For our parameters these are $\alpha_A=11/29$, $\beta_h=3/29$, $\gamma_H=39/29$, $\gamma_d=53/58$,
$a_0=41/58$ and $a_1=76/29$.

\subsection{A quadrature formula for the coefficient of $G$}\label{ssec:quadrature}
Let $\varepsilon=x-G$. Then $\Delta_N(x)=\det(\cA_N+\varepsilon\cL_N)$. The question is how large the columns
of $\cL_N$ are compared with those of $\cA_N$. By \eqref{eq:L} the entries of $\cL_N$ are linear
functionals of the coefficients of the row polynomials. These polynomials have bounded degree, so the
functionals can be written exactly as finite combinations of point evaluations. Let $t_1,\ldots,t_J$ be
distinct points, with $J$ at least the number of degrees of freedom of the rows. Lagrange interpolation
gives
\begin{equation}\label{eq:quad}
 \cL_N(r,k)=\sum_{a=1}^JA_r(t_a)\,\beta_{a,k},
\end{equation}
where the weight $\beta_{a,k}$ is the value of the same functional on the $a$-th Lagrange basis polynomial
$\ell_a$. This is a quadrature formula of Gauss type. It makes the columns of $\cL_N$ combinations of point
values of the row functions that define $\cA_N$. So the perturbed determinant can be treated by the energy
method. The price is the size of the weights.

We take $J=d_0N=53N$ nodes. Put $g_*=401/2500$, $a_*=2g_*/(1+g_*^2)=2005000/6410801$ and $J'=\lceil J/2\rceil$.
The nodes are the points
\[
 t_{j,\pm}=\pm\sqrt{a_*^2+(1-a_*^2)(1+s_j)/2},\qquad T_{J'}(s_j)=0,
\]
that is, the roots of $T_{J'}\bigl((2t^2-(1+a_*^2))/(1-a_*^2)\bigr)$. We keep all $J'$ positive roots and the
first $J-J'$ negative ones, so for odd $J$ exactly one root is dropped. The nodes lie in
$S=(-1,-a_*)\cup(a_*,1)$, and in the coordinate $x$ with $t=2x/(1+x^2)$ every node has $|x|>g_*$. The nodes
avoid a neighbourhood of $x=0$, where the fields of the energy method are largest, and they are Chebyshev
nodes for $S$. Dropping a root changes the nodal polynomial by a linear factor, so it changes the
derivative bound and the size of the Lagrange polynomials only by factors polynomial in $N$. So an odd
number of nodes does not change the exponent. This is needed, since $J=53N$ is odd for odd $N$.

Let $g_S$ be the Green function of $\C\setminus\bar S$ with pole at infinity. In closed form
\[
 g_S(z)=\tfrac12\operatorname{arccosh}\Bigl(\frac{|z^2-1|+|z^2-a_*^2|}{1-a_*^2}\Bigr).
\]
The Lagrange polynomials satisfy $|\ell_a(z)|\le e^{53Ng_S(z)+o(N)}$ at distance at least $c/N$ from the
nodes, by the derivative identity for Chebyshev polynomials at their roots. The weight of column $k$ is a
coefficient functional of $\ell_a$ multiplied by the column function. After removal of the common factor of
the rows it equals
\[
 \beta'_k=4\,[t^0]\,F_k(t)\,(1-t^{-2})^{-1/2},\qquad F_k(t)=t^{\,A-b-q-k}(1-t)^h(t-1)^q\ell_a(t),
\]
where $(1-t^{-2})^{-1/2}=\sum_lc_lt^{-2l}$ with $c_l=\binom{2l}l4^{-l}$. The next lemma bounds it by a contour
that depends on $k$.

\begin{lemma}[keyhole bound]\label{lem:keyhole}
Let $P$ be a real polynomial, $E\ge0$, and $F(t)=t^{-E}P(t)$. Let $0<r\le1$, and suppose that $|F|\le M$ on
the circle $|t|=r$ and $|F|\le M'$ on the segments $r\le|t|\le1$ of the real axis. Then
\[
 \Bigl|4\sum_{l\ge0}c_l\,[t^{E+2l}]P\Bigr|\le4(1+E)\max(M,M').
\]
\end{lemma}

\begin{proof}
Write $F=\Pi+R$ with $\Pi$ a polynomial and $R=\sum_{i=1}^E\pi_it^{-i}$ the principal part at $0$. The
functional sees only $\Pi$. Since $c_l=\frac1\pi\int_{-1}^1t^{2l}(1-t^2)^{-1/2}dt$, its value is
$\frac4\pi\int_{-1}^1\Pi(t)(1-t^2)^{-1/2}dt$, which is at most $4\max_{[-1,1]}|\Pi|$. By Cauchy's estimate
$|\pi_i|\le Mr^i$, so $|R(t)|\le EM$ for $r\le|t|\le1$, and $|\Pi|\le M'+EM$ there. For $|t|<r$ the maximum
principle gives $|\Pi(t)|\le\max_{|t'|=r}|\Pi|\le(1+E)M$.
\end{proof}

The formal proof replaces the integral by the finite Gauss--Chebyshev identity
$\frac1M\sum_{j<M}\cos^{2l}\frac{(2j+1)\pi}{2M}=c_l$ for $l<M$, and Cauchy's estimate by the finite Cauchy formula
at roots of unity. So no integration is needed. For $0\le y\le1$ and $z\in\C$ let
\[
 \Phi_y(z)=(10-58y)\log|z|+6\log|1-z|+53g_S(z).
\]
Here $10=(A-b-q)/N$ and $6=(q+h)/N$. With $y=k/n$, Lemma~\ref{lem:keyhole} gives
\begin{equation}\label{eq:wk}
 |\beta'_k|\le e^{(w(y)+o(1))N},\qquad w(y)=\min_r\max_{z\in\Gamma_r}\Phi_y(z),
\end{equation}
where $\Gamma_r$ is the circle $|z|=r$ together with the segments $r\le|z|\le1$ of the real axis. For $y$ near
$1$ the factor $|z|^{10-58y}$ favours a large circle, and for small $y$ a small one. The circle $|z|=1.001$ alone
gives $w(y)\le52.79$ for every $y$. The keyhole contour improves this only slightly at the worst column. Its
use is elsewhere: it lowers $w(y)$ at the columns where the binomial factor of Lemma~\ref{lem:binomvand} is
largest.

\subsection{Column expansion and Andr\'eief's identity}\label{ssec:andreief}
Write $\mathcal P(r,a)=A_r(t_a)$. Expanding $\det(\cA_N+\varepsilon\mathcal P\beta)$ multilinearly in the
columns gives a sum over maps $f$ from the columns to $\{\varnothing,1,\ldots,J\}$ of
$\prod_{f(k)\ne\varnothing}\varepsilon\beta_{f(k),k}\det M_f$. Here $M_f$ is the matrix whose $k$-th column is
the $k$-th column of $\cA_N$ if $f(k)=\varnothing$, and the vector $(A_r(t_{f(k)}))_r$ otherwise.

\begin{lemma}\label{lem:expansion}
Suppose that $|\det M_f|\le M_0\prod_{f(k)\ne\varnothing}C_k$ for every $f$, and that $|\beta_{a,k}|\le B_k$.
Then
\[
 \bigl|\det(\cA_N+\varepsilon\mathcal P\beta)-\det\cA_N\bigr|\le M_0\Bigl(\prod_k(1+J|\varepsilon|B_kC_k)-1\Bigr).
\]
\end{lemma}

The lemma allows a different constant for each column. Each $M_f$ has $n'=n-m$ columns that are integrals
and $m$ columns that are point values. Andr\'eief's identity extends to this situation.

\begin{lemma}\label{lem:andreief}
Let $m_1,m_2$ be finite measures, and let $F_r$ $(r<n'+m)$, $F'_i$ $(i<n')$ and $k$ be integrable functions.
Let $B$ be the matrix with $B(r,i)=\iint F_r(t)k(t,s)F'_i(s)\,dm_1(t)\,dm_2(s)$ for $i<n'$ and
$B(r,n'+j)=F_r(t'_j)$ for $j<m$. Then
\begin{multline*}
 (n'!)^2\det B=\int\!\!\int\det\bigl[F_r((t\Vert t')_c)\bigr]_{r,c}\\
 \cdot\det[k(t_i,s_j)]_{i,j<n'}\det[F'_i(s_j)]_{j,i<n'}\,dm_1^{n'}(t)\,dm_2^{n'}(s),
\end{multline*}
where $t\Vert t'$ is the list of the $n'$ variables followed by the $m$ prescribed points.
\end{lemma}

\begin{proof}
For $m=0$ this is Andr\'eief's identity applied twice \cite{Forrester2018Andreief}. In general, expand
$\det B$ multilinearly in its first $n'$ columns. This introduces one pair of variables $(t_i,s_i)$ for each
of them. Then antisymmetrize, first in the $t$ variables and then in the $s$ variables. The last $m$ columns
are constant throughout.
\end{proof}

With $k(t,s)=1/(1-ts)$ the middle determinant is a Cauchy determinant. The third is a generalized
Vandermonde determinant $\prod_js_j^b(1-s_j)^q\det[s_j^{E_i}]$. Here the exponents $E_i<n$ are those of the
columns that remain integrals, and the missing exponents are the replaced columns. We bound the generalized
Vandermonde determinant by a product over the missing exponents.

\begin{lemma}\label{lem:binomvand}
Let $s_1,\ldots,s_{n'}\in[0,1]$, let $E_1<\dots<E_{n'}<n$ be integers, and let $\{k_1,\ldots,k_m\}$ be the
complement of $\{E_i\}$ in $\{0,\ldots,n-1\}$. Then
\[
 \bigl|\det[s_l^{E_i}]\bigr|\le\prod_{j=1}^m\mathcal B(k_j)\prod_{i<l}|s_i-s_l|,\qquad
 \mathcal B(k)=\sum_{i=0}^{m-1}\binom{n'}{n'-k+i},
\]
with $\binom{n'}r=0$ for $r\notin[0,n']$.
\end{lemma}

\begin{proof}
Let $P(X)=\prod_l(X-s_l)=\sum_r(-1)^re_r(s)X^{n'-r}$. Multiplying the matrix by a unitriangular matrix built
from $P$ gives $|\det[s_l^{E_i}]|=\prod_{i<l}|s_i-s_l|\cdot|\det Q|$, where $Q$ is the $m\times m$ matrix with
entries $Q_{j,i}=[X^{k_j-i}]P$. So $|Q_{j,i}|=e_{n'-k_j+i}(s)\le\binom{n'}{n'-k_j+i}$, since $s\in[0,1]^{n'}$.
Hadamard's inequality, applied to the rows of $Q$ in the $\ell^1$ norm, gives the claim.
\end{proof}

\subsection{The weight of one column}\label{ssec:weights}
The bound of Lemma~\ref{lem:binomvand} does not depend on the points $s_l$. Since $\binom{n'}r\le\binom nr$,
\[
 \log\mathcal B(k)\le n_0\tilde h_2(k/n)N+O(\log N),
\]
where $\tilde h_2(y)$ is the binary entropy $h_2(y)=-y\log y-(1-y)\log(1-y)$ for $y\ge\frac12$ and $\log2$ for
$y<\frac12$. So the replaced column $k$ contributes at most $e^{(w(y)+n_0\tilde h_2(y)+o(1))N}$, with $y=k/n$.
Put
\[
 W_0\ge\sup_{0\le y\le1}\bigl(w(y)+58\tilde h_2(y)\bigr).
\]
Appendix~\ref{app:field} proves $W_0=\WZero$.

The two factors peak at different columns. The weight $w(y)$ increases up to $y\approx0.84$, while
$\tilde h_2(y)$ is largest on $[0,\frac12]$. So the maximum of $w+58\tilde h_2$ is about $87.96$, at
$y\approx0.62$, well below the sum $52.79+40.20$ of the separate maxima (Figure~\ref{fig:columns}). This saving of
about $5$ in $\rho$ is the reason for treating the columns one at a time.

\begin{figure}[ht]
\centering
\setlength{\unitlength}{1cm}
\begin{picture}(12.6,5.6)(-1.0,-0.8)
\put(0,0){\line(1,0){10.2}}
\put(0,0){\line(0,1){4.9}}
\put(0.000,0){\line(0,-1){0.12}}
\put(0.000,-0.4){\makebox(0,0){$0$}}
\put(2.000,0){\line(0,-1){0.12}}
\put(2.000,-0.4){\makebox(0,0){$0.2$}}
\put(4.000,0){\line(0,-1){0.12}}
\put(4.000,-0.4){\makebox(0,0){$0.4$}}
\put(6.000,0){\line(0,-1){0.12}}
\put(6.000,-0.4){\makebox(0,0){$0.6$}}
\put(8.000,0){\line(0,-1){0.12}}
\put(8.000,-0.4){\makebox(0,0){$0.8$}}
\put(10.000,0){\line(0,-1){0.12}}
\put(10.000,-0.4){\makebox(0,0){$1$}}
\put(0,0.000){\line(-1,0){0.12}}
\put(-0.25,0.000){\makebox(0,0)[r]{$0$}}
\put(0,1.000){\line(-1,0){0.12}}
\put(-0.25,1.000){\makebox(0,0)[r]{$20$}}
\put(0,2.000){\line(-1,0){0.12}}
\put(-0.25,2.000){\makebox(0,0)[r]{$40$}}
\put(0,3.000){\line(-1,0){0.12}}
\put(-0.25,3.000){\makebox(0,0)[r]{$60$}}
\put(0,4.000){\line(-1,0){0.12}}
\put(-0.25,4.000){\makebox(0,0)[r]{$80$}}
\put(10.15,-0.4){\makebox(0,0)[l]{$y$}}
\qbezier(0.000,0.586)(0.100,0.621)(0.200,0.656)
\qbezier(0.200,0.656)(0.300,0.691)(0.400,0.725)
\qbezier(0.400,0.725)(0.500,0.760)(0.600,0.795)
\qbezier(0.600,0.795)(0.700,0.830)(0.800,0.865)
\qbezier(0.800,0.865)(0.900,0.900)(1.000,0.935)
\qbezier(1.000,0.935)(1.100,0.970)(1.200,1.005)
\qbezier(1.200,1.005)(1.300,1.040)(1.400,1.075)
\qbezier(1.400,1.075)(1.500,1.110)(1.600,1.144)
\qbezier(1.600,1.144)(1.700,1.179)(1.800,1.214)
\qbezier(1.800,1.214)(1.900,1.249)(2.000,1.284)
\qbezier(2.000,1.284)(2.100,1.319)(2.200,1.354)
\qbezier(2.200,1.354)(2.300,1.389)(2.400,1.424)
\qbezier(2.400,1.424)(2.500,1.459)(2.600,1.494)
\qbezier(2.600,1.494)(2.700,1.529)(2.800,1.563)
\qbezier(2.800,1.563)(2.900,1.598)(3.000,1.633)
\qbezier(3.000,1.633)(3.100,1.668)(3.200,1.703)
\qbezier(3.200,1.703)(3.300,1.738)(3.400,1.773)
\qbezier(3.400,1.773)(3.500,1.808)(3.600,1.843)
\qbezier(3.600,1.843)(3.700,1.877)(3.800,1.911)
\qbezier(3.800,1.911)(3.900,1.944)(4.000,1.976)
\qbezier(4.000,1.976)(4.100,2.006)(4.200,2.037)
\qbezier(4.200,2.037)(4.300,2.065)(4.400,2.094)
\qbezier(4.400,2.094)(4.500,2.121)(4.600,2.148)
\qbezier(4.600,2.148)(4.700,2.174)(4.800,2.199)
\qbezier(4.800,2.199)(4.900,2.223)(5.000,2.247)
\qbezier(5.000,2.247)(5.100,2.269)(5.200,2.291)
\qbezier(5.200,2.291)(5.300,2.312)(5.400,2.333)
\qbezier(5.400,2.333)(5.500,2.353)(5.600,2.372)
\qbezier(5.600,2.372)(5.700,2.390)(5.800,2.408)
\qbezier(5.800,2.408)(5.900,2.425)(6.000,2.442)
\qbezier(6.000,2.442)(6.100,2.457)(6.200,2.472)
\qbezier(6.200,2.472)(6.300,2.486)(6.400,2.501)
\qbezier(6.400,2.501)(6.500,2.513)(6.600,2.526)
\qbezier(6.600,2.526)(6.700,2.537)(6.800,2.549)
\qbezier(6.800,2.549)(6.900,2.559)(7.000,2.569)
\qbezier(7.000,2.569)(7.100,2.578)(7.200,2.587)
\qbezier(7.200,2.587)(7.300,2.594)(7.400,2.602)
\qbezier(7.400,2.602)(7.500,2.608)(7.600,2.614)
\qbezier(7.600,2.614)(7.700,2.619)(7.800,2.624)
\qbezier(7.800,2.624)(7.900,2.628)(8.000,2.631)
\qbezier(8.000,2.631)(8.100,2.633)(8.200,2.636)
\qbezier(8.200,2.636)(8.300,2.636)(8.400,2.637)
\qbezier(8.400,2.637)(8.500,2.637)(8.600,2.637)
\qbezier(8.600,2.637)(8.700,2.637)(8.800,2.637)
\qbezier(8.800,2.637)(8.900,2.637)(9.000,2.637)
\qbezier(9.000,2.637)(9.100,2.637)(9.200,2.637)
\qbezier(9.200,2.637)(9.300,2.637)(9.400,2.637)
\qbezier(9.400,2.637)(9.500,2.637)(9.600,2.637)
\qbezier(9.600,2.637)(9.700,2.637)(9.800,2.637)
\qbezier(9.800,2.637)(9.900,2.637)(10.000,2.637)
\qbezier(0.000,2.010)(0.100,2.010)(0.200,2.010)
\qbezier(0.200,2.010)(0.300,2.010)(0.400,2.010)
\qbezier(0.400,2.010)(0.500,2.010)(0.600,2.010)
\qbezier(0.600,2.010)(0.700,2.010)(0.800,2.010)
\qbezier(0.800,2.010)(0.900,2.010)(1.000,2.010)
\qbezier(1.000,2.010)(1.100,2.010)(1.200,2.010)
\qbezier(1.200,2.010)(1.300,2.010)(1.400,2.010)
\qbezier(1.400,2.010)(1.500,2.010)(1.600,2.010)
\qbezier(1.600,2.010)(1.700,2.010)(1.800,2.010)
\qbezier(1.800,2.010)(1.900,2.010)(2.000,2.010)
\qbezier(2.000,2.010)(2.100,2.010)(2.200,2.010)
\qbezier(2.200,2.010)(2.300,2.010)(2.400,2.010)
\qbezier(2.400,2.010)(2.500,2.010)(2.600,2.010)
\qbezier(2.600,2.010)(2.700,2.010)(2.800,2.010)
\qbezier(2.800,2.010)(2.900,2.010)(3.000,2.010)
\qbezier(3.000,2.010)(3.100,2.010)(3.200,2.010)
\qbezier(3.200,2.010)(3.300,2.010)(3.400,2.010)
\qbezier(3.400,2.010)(3.500,2.010)(3.600,2.010)
\qbezier(3.600,2.010)(3.700,2.010)(3.800,2.010)
\qbezier(3.800,2.010)(3.900,2.010)(4.000,2.010)
\qbezier(4.000,2.010)(4.100,2.010)(4.200,2.010)
\qbezier(4.200,2.010)(4.300,2.010)(4.400,2.010)
\qbezier(4.400,2.010)(4.500,2.010)(4.600,2.010)
\qbezier(4.600,2.010)(4.700,2.010)(4.800,2.010)
\qbezier(4.800,2.010)(4.900,2.010)(5.000,2.010)
\qbezier(5.000,2.010)(5.100,2.009)(5.200,2.008)
\qbezier(5.200,2.008)(5.300,2.004)(5.400,2.001)
\qbezier(5.400,2.001)(5.500,1.995)(5.600,1.989)
\qbezier(5.600,1.989)(5.700,1.981)(5.800,1.973)
\qbezier(5.800,1.973)(5.900,1.962)(6.000,1.952)
\qbezier(6.000,1.952)(6.100,1.939)(6.200,1.926)
\qbezier(6.200,1.926)(6.300,1.910)(6.400,1.895)
\qbezier(6.400,1.895)(6.500,1.877)(6.600,1.859)
\qbezier(6.600,1.859)(6.700,1.838)(6.800,1.818)
\qbezier(6.800,1.818)(6.900,1.795)(7.000,1.772)
\qbezier(7.000,1.772)(7.100,1.746)(7.200,1.720)
\qbezier(7.200,1.720)(7.300,1.691)(7.400,1.662)
\qbezier(7.400,1.662)(7.500,1.630)(7.600,1.598)
\qbezier(7.600,1.598)(7.700,1.563)(7.800,1.528)
\qbezier(7.800,1.528)(7.900,1.490)(8.000,1.451)
\qbezier(8.000,1.451)(8.100,1.409)(8.200,1.367)
\qbezier(8.200,1.367)(8.300,1.321)(8.400,1.275)
\qbezier(8.400,1.275)(8.500,1.225)(8.600,1.174)
\qbezier(8.600,1.174)(8.700,1.119)(8.800,1.064)
\qbezier(8.800,1.064)(8.900,1.003)(9.000,0.943)
\qbezier(9.000,0.943)(9.100,0.876)(9.200,0.808)
\qbezier(9.200,0.808)(9.300,0.733)(9.400,0.658)
\qbezier(9.400,0.658)(9.500,0.573)(9.600,0.487)
\qbezier(9.600,0.487)(9.700,0.386)(9.800,0.284)
\qbezier(9.800,0.284)(9.900,0.142)(10.000,0.000)
\qbezier(0.000,2.596)(0.100,2.631)(0.200,2.666)
\qbezier(0.200,2.666)(0.300,2.701)(0.400,2.736)
\qbezier(0.400,2.736)(0.500,2.771)(0.600,2.805)
\qbezier(0.600,2.805)(0.700,2.840)(0.800,2.875)
\qbezier(0.800,2.875)(0.900,2.910)(1.000,2.945)
\qbezier(1.000,2.945)(1.100,2.980)(1.200,3.015)
\qbezier(1.200,3.015)(1.300,3.050)(1.400,3.085)
\qbezier(1.400,3.085)(1.500,3.120)(1.600,3.155)
\qbezier(1.600,3.155)(1.700,3.190)(1.800,3.224)
\qbezier(1.800,3.224)(1.900,3.259)(2.000,3.294)
\qbezier(2.000,3.294)(2.100,3.329)(2.200,3.364)
\qbezier(2.200,3.364)(2.300,3.399)(2.400,3.434)
\qbezier(2.400,3.434)(2.500,3.469)(2.600,3.504)
\qbezier(2.600,3.504)(2.700,3.539)(2.800,3.574)
\qbezier(2.800,3.574)(2.900,3.608)(3.000,3.643)
\qbezier(3.000,3.643)(3.100,3.678)(3.200,3.713)
\qbezier(3.200,3.713)(3.300,3.748)(3.400,3.783)
\qbezier(3.400,3.783)(3.500,3.818)(3.600,3.853)
\qbezier(3.600,3.853)(3.700,3.887)(3.800,3.921)
\qbezier(3.800,3.921)(3.900,3.954)(4.000,3.986)
\qbezier(4.000,3.986)(4.100,4.016)(4.200,4.047)
\qbezier(4.200,4.047)(4.300,4.076)(4.400,4.104)
\qbezier(4.400,4.104)(4.500,4.131)(4.600,4.158)
\qbezier(4.600,4.158)(4.700,4.184)(4.800,4.209)
\qbezier(4.800,4.209)(4.900,4.233)(5.000,4.257)
\qbezier(5.000,4.257)(5.100,4.278)(5.200,4.299)
\qbezier(5.200,4.299)(5.300,4.317)(5.400,4.334)
\qbezier(5.400,4.334)(5.500,4.348)(5.600,4.361)
\qbezier(5.600,4.361)(5.700,4.371)(5.800,4.381)
\qbezier(5.800,4.381)(5.900,4.387)(6.000,4.393)
\qbezier(6.000,4.393)(6.100,4.396)(6.200,4.398)
\qbezier(6.200,4.398)(6.300,4.397)(6.400,4.395)
\qbezier(6.400,4.395)(6.500,4.390)(6.600,4.385)
\qbezier(6.600,4.385)(6.700,4.376)(6.800,4.367)
\qbezier(6.800,4.367)(6.900,4.354)(7.000,4.341)
\qbezier(7.000,4.341)(7.100,4.323)(7.200,4.306)
\qbezier(7.200,4.306)(7.300,4.285)(7.400,4.264)
\qbezier(7.400,4.264)(7.500,4.238)(7.600,4.212)
\qbezier(7.600,4.212)(7.700,4.182)(7.800,4.152)
\qbezier(7.800,4.152)(7.900,4.117)(8.000,4.082)
\qbezier(8.000,4.082)(8.100,4.042)(8.200,4.003)
\qbezier(8.200,4.003)(8.300,3.957)(8.400,3.912)
\qbezier(8.400,3.912)(8.500,3.862)(8.600,3.812)
\qbezier(8.600,3.812)(8.700,3.756)(8.800,3.701)
\qbezier(8.800,3.701)(8.900,3.641)(9.000,3.580)
\qbezier(9.000,3.580)(9.100,3.513)(9.200,3.446)
\qbezier(9.200,3.446)(9.300,3.370)(9.400,3.295)
\qbezier(9.400,3.295)(9.500,3.210)(9.600,3.124)
\qbezier(9.600,3.124)(9.700,3.023)(9.800,2.921)
\qbezier(9.800,2.921)(9.900,2.779)(10.000,2.637)
\multiput(0,4.650)(0.2,0){51}{\line(1,0){0.08}}
\put(0.2,4.870){\makebox(0,0)[l]{$52.79+40.20$}}
\put(6.200,4.398){\makebox(0,0){$\bullet$}}
\put(6.200,4.048){\makebox(0,0){$87.96$}}
\put(10.15,2.637){\makebox(0,0)[l]{$w(y)$}}
\put(2.0,2.260){\makebox(0,0)[l]{$58\tilde h_2(y)$}}
\put(2.6,3.804){\makebox(0,0)[r]{sum}}
\end{picture}

\caption{The weight $w(y)$ of column $k=yn$, the binomial factor $58\tilde h_2(y)$, and their sum, per unit
of $N$. The values were computed numerically on a grid of step $0.02$ and are not part of the proof. The
certified bound for the sum is $W_0=\WZero$. The dotted line is the sum of the separate maxima.}
\label{fig:columns}
\end{figure}

The remaining factor of $\rho$ is $n_0c_2$, the increase of the energy bound caused by the prescribed
columns. It comes from Section~\ref{sec:energy}, and with it $\rho=58\cdot\Ctwo+\WZero=\Rho$.
